\chapter{Model implementation}
\label{app:model_impl}

\begin{lstlisting}[style=Rcode, caption={Model implementation. This code was used to generate the SWE maps.}, label={lst:swe}]
###################################################################
#                    SWE model implementation                     #
###################################################################
rm(list = ls()); invisible(gc())

library(terra)
library(ncdf4)

years <- 1951:2023
fname_in <- "input.dat"
setwd("/home/filippo/Desktop/Codicini/Master_Thesis/SWE_model/")



# Fun. to compute solid precipitations through a single temperature
# threshold approach: the threshold, together with the daily min 
# and max temperatures, determines the fraction of precipitation 
# that falls as snow or that falls as rain.
compute_solid_precipitations <- function(prec, tmnd, tmxd, t_th){
  
  precs <- prec 
  
  # Threshold temperature is comprised within daily tmin and tmax
  mask <- tmnd <= t_th & tmxd > t_th & !is.na(prec) & !is.na(tmnd) & !is.na(tmxd) & prec > 0
  precs[mask] <- prec[mask]*(t_th - tmnd[mask])/(tmxd[mask] - tmnd[mask])
  rm(mask); invisible(gc())
  

  # Daily temperatures are too high for snowfall to occur
  mask <- tmnd > t_th & !is.na(prec) & !is.na(tmnd) & !is.na(tmxd) & prec > 0
  precs[mask] <- 0
  rm(mask); invisible(gc())
  
  
  # Alligning precipitation and temperature masks, in order to 
  # compute SWE only on the overlap area
  mask <- !is.na(precs) & (is.na(tmnd) | is.na(tmxd))
  precs[mask] <- NA
  rm(mask); invisible(gc())
  
  return(precs)
}


# Fun. to compute oscillating DDF. The period of the oscillation is 
# 366 days, but for non-leap years the 29 February value is deleted
compute_ddf <- function(year, ddf_ave, ddf_ampl){
  
  # DDF computation
  len <- 366; offset <- 81;idx <- 1:len
  ddf <- ddf_ave + ddf_ampl * sin(2 * pi * (idx - offset)/len)
  
  # Leap on non-leap year?
  mask <- year %% 4 == 0 && (year %% 100 != 0 || year %% 400 == 0)
  if(!mask) ddf <- ddf[-60]
  return(ddf)
}


# Fun. to compute snowmelt based on Magnusson approach. The degree 
# day is corrected to ensure smoothness at threshold temperature.
compute_melt <- function(year, tmean, ddf_ave, ddf_ampl, expfact){
  
  # Computing oscillating DDF and degree day
  ddf <- compute_ddf(year, ddf_ave, ddf_ampl)
  deg_day <- tmean + expfact * log(1 + exp(-tmean/expfact))
  
  # Snowmelt computation (dimension match?)
  stopifnot(dim(tmean)[3] == length(ddf))
  melt <- sweep(deg_day, 3, ddf, "*")
  
  # Looking for eventual overflows
  stopifnot(all(is.finite(melt) | is.na(melt)))
  return(melt)
}


# Function to create the container that stores the SWE map, updated 
# on a daily basis across the 1951--2023 period.
create_swe_container <- function(fname){
  nc <- nc_open(fname)
  prec <- ncvar_get(nc, "total_precipitation")
  nc_close(nc)
  
  appo <- array(0, dim = c(dim(prec)[1], dim(prec)[2]))
  rm(prec); invisible(gc())
  return(appo)
}


# Function to save swe maps during a given solar year to netCDF. 
# The structure of the input files is used in order to keep the 
# same architecture for input and output files
save_annual_swe <- function(total, fname_template, fname_out){
  
  nc <- nc_open(fname_template)
  v <- nc$var[["total_precipitation"]]
  if(is.null(v)){
    on.exit(nc_close(nc))
    stop("Variable 'total_precipitation' not found. You can 
      choose between: ", paste(names(nc$var), collapse = ", "))
  }
  stopifnot(length(v$dim) == 3)
  
  dims <- vector("list", 3)
  for(i in seq_along(v$dim)){
    d <- v$dim[[i]]
    dims[[i]] <- ncdim_def(
      name = d$name, units = d$units, vals = d$vals,
      unlim = d$unlim, create_dimvar = isTRUE(d$create_dimvar),
      calendar = if(is.null(d$calendar)) NA else d$calendar)
  }
  nc_close(nc); rm(nc, v); invisible(gc())
  
  
  # Checking if netCDF container and SWE container dimensions match
  lens <- vapply(dims, function(d) as.integer(d$len), integer(1))
  stopifnot(identical(lens, as.integer(dim(total))))
  
  swe <- ncvar_def(
    name = "swe", units = "mm", dim = dims, missval = -9999, 
    longname = "Snow water equivalent (SWE)", prec = "float", 
    compression = 5)
  
  nc_out <- nc_create(fname_out, swe, force_v4 = TRUE)
  ncvar_put(nc_out, "swe", total)
  nc_close(nc_out)
  
  rm(nc_out, swe, dims); invisible(gc())
}










# Importing model parameters
df_in <- read.table(fname_in, header = TRUE)
ddf_ave <- as.numeric(df_in$ddf_ave)
ddf_ampl <- as.numeric(df_in$ddf_ampl)
expfact <- as.numeric(df_in$expfact)
t_th <- as.numeric(df_in$tlim)

if((ddf_ave - ddf_ampl) < 0){ stop("DDF not correct: stopping! ") }

appo <- create_swe_container("Dataset/PCPD/1951.nc")
for(y in years){
  
  # Importing temperature and precipitation files
  fname_prec <- paste0("Dataset/PCPD/", y, ".nc")
  fname_temp <- paste0("Dataset/TEMP/", y, ".nc")
  
  prec <- rast(fname_prec)
  temp <- rast(fname_temp)
  
  if(!compareGeom(prec, temp, stopOnError = FALSE)){
    stop(paste0("No compatible input grids during ", y))
  }
  rm(prec, temp)
  invisible(gc())
  
  
  
  # Loading precipitation and temperature grids 
  # (every layer is one daily map)
  nc <- nc_open(fname_prec)
  prec <- ncvar_get(nc, "total_precipitation")
  nc_close(nc)
  
  nc <- nc_open(fname_temp)
  tmnd <- ncvar_get(nc, "tmnd")
  tmxd <- ncvar_get(nc, "tmxd")
  stopifnot(identical(dim(tmnd), dim(tmxd)))
  tmean <- (tmnd + tmxd)/2
  nc_close(nc)
  
  rm(nc); invisible(gc());
  stopifnot(identical(dim(tmean), dim(prec)))
  
  
  
  
  
  # Computing solid precipitations
  precs <- compute_solid_precipitations(prec, tmnd, tmxd, t_th)
  rm(prec, tmnd, tmxd)
  invisible(gc())
  
  
  # Computing annual melt
  melt <- compute_melt(y, tmean, ddf_ave, ddf_ampl, expfact)
  rm(tmean); invisible(gc())
  
  total <- array(NA_real_, dim = dim(melt))
  for(i in 1:dim(total)[3]){
    
    # Adding solid precipitations and melt to the SWE map. If some 
    # negative values appear, it means that snow has melted all the 
    # way to the ground (these pixels must be set to zero)
    appo <- appo + precs[, , i] - melt[, , i]
    
    mask <- appo < 0 & !is.na(appo)
    appo[mask] <- 0
    
    total[, , i] <- round(appo, 1)
  }
  rm(mask); invisible(gc())
  

  # Saving SWE data to file
  total[is.na(total)] <- -9999
  fname_template <- paste0("Dataset/PCPD/", y, ".nc")
  fname_out <- paste0("Results/SWE_", y, ".nc")
  save_annual_swe(total, fname_template, fname_out)
  
  rm(precs, melt, total); invisible(gc())
  cat("Computed SWE maps for: ", y, "       Number of pixel NAs: ", sum(is.na(appo)),"\n")
}
\end{lstlisting}